\documentclass[12pt,letterpaper]{article}
\usepackage{graphicx,textcomp}
\usepackage{natbib}
\usepackage{setspace}
\usepackage{fullpage}
\usepackage{color}
\usepackage[reqno]{amsmath}
\usepackage{amsthm}
\usepackage{fancyvrb}
\usepackage{amssymb,enumerate}
\usepackage[all]{xy}
\usepackage{endnotes}
\usepackage{lscape}
\newtheorem{com}{Comment}
\usepackage{float}
\usepackage{hyperref}
\newtheorem{lem} {Lemma}
\newtheorem{prop}{Proposition}
\newtheorem{thm}{Theorem}
\newtheorem{defn}{Definition}
\newtheorem{cor}{Corollary}
\newtheorem{obs}{Observation}
\usepackage[compact]{titlesec}
\usepackage{dcolumn}
\usepackage{tikz}
\usetikzlibrary{arrows}
\usepackage{multirow}
\usepackage{xcolor}
\newcolumntype{.}{D{.}{.}{-1}}
\newcolumntype{d}[1]{D{.}{.}{#1}}
\definecolor{light-gray}{gray}{0.65}
\usepackage{url}
\usepackage{listings}
\usepackage{color}

\definecolor{codegreen}{rgb}{0,0.6,0}
\definecolor{codegray}{rgb}{0.5,0.5,0.5}
\definecolor{codepurple}{rgb}{0.58,0,0.82}
\definecolor{backcolour}{rgb}{0.95,0.95,0.92}

\lstdefinestyle{mystyle}{
	backgroundcolor=\color{backcolour},   
	commentstyle=\color{codegreen},
	keywordstyle=\color{magenta},
	numberstyle=\tiny\color{codegray},
	stringstyle=\color{codepurple},
	basicstyle=\footnotesize,
	breakatwhitespace=false,         
	breaklines=true,                 
	captionpos=b,                    
	keepspaces=true,                 
	numbers=left,                    
	numbersep=5pt,                  
	showspaces=false,                
	showstringspaces=false,
	showtabs=false,                  
	tabsize=2
}
\lstset{style=mystyle}
\newcommand{\Sref}[1]{Section~\ref{#1}}
\newtheorem{hyp}{Hypothesis}


\title{Problem Set 4}
\date{Ella Lang}
\author{Applied Stats/Quant Methods 1}


\begin{document}
	\maketitle
	\section*{Instructions}
	\begin{itemize}
		\item Please show your work! You may lose points by simply writing in the answer. If the problem requires you to execute commands in \texttt{R}, please include the code you used to get your answers. Please also include the \texttt{.R} file that contains your code. If you are not sure if work needs to be shown for a particular problem, please ask.
		\item Your homework should be submitted electronically on GitHub.
		\item This problem set is due before 23:59 on Thursday November 26, 2026. No late assignments will be accepted.
	\end{itemize}



	\vspace{.5cm}
\section*{Question 1: Economics}
\vspace{.25cm}
\noindent 	
In this question, use the \texttt{prestige} dataset in the \texttt{car} library. First, run the following commands:

\begin{verbatim}
install.packages(car)
library(car)
data(Prestige)
help(Prestige)
\end{verbatim} 


\noindent We would like to study whether individuals with higher levels of income have more prestigious jobs. Moreover, we would like to study whether professionals have more prestigious jobs than blue and white collar workers.

\newpage
\begin{enumerate}
	
	\item [(a)]
	Create a new variable \texttt{professional} by recoding the variable \texttt{type} so that professionals are coded as $1$, and blue and white collar workers are coded as $0$ (Hint: \texttt{ifelse}).
	
	\noindent I first checked the variable using levels(), class() and table(). Linear regression only includes complete rows in the calculation so $<$NA$>$'s did not need to be removed. I then used the ifelse function to create the Professional variable with professionals coded as 1 and blue and white collar workers coded as 0. I checked to see if the function worked using with the table() function.  \\
	\lstinputlisting[language=R, firstline=45,lastline=48]{PS04.my.answers.el.la.R}
	\begin{lstlisting}[language=R]
	"bc"   "prof" "wc"  
	"factor"
	bc prof   wc 
	44   31   23 	\end{lstlisting}
	\lstinputlisting[language=R, firstline=50,lastline=51]{PS04.my.answers.el.la.R}
		\begin{lstlisting}[language=R]
		0  1 
		67 31  	\end{lstlisting}
	
	\item [(b)]
	Run a linear model with \texttt{prestige} as an outcome and \texttt{income}, \texttt{professional}, and the interaction of the two as predictors (Note: this is a continuous $\times$ dummy interaction.)
	
	\noindent I ran a linear regression with the lm function, and using the * symbol to include the interaction of the two predictors. The output of the regression is on the following page. \\
	\lstinputlisting[language=R, firstline=56,lastline=57]{PS04.my.answers.el.la.R}
	\begin{table}[!htbp] \centering 
		\caption{} 
		\label{} 
		\begin{tabular}{@{\extracolsep{5pt}}lc} 
			\\[-1.8ex]\hline 
			\hline \\[-1.8ex] 
			& \multicolumn{1}{c}{\textit{Dependent variable:}} \\ 
			\cline{2-2} 
			\\[-1.8ex] & Prestige \\ 
			\hline \\[-1.8ex] 
			Income & 0.003$^{***}$ \\ 
			& (0.0005) \\ 
			& \\ 
			Professional & 37.781$^{***}$ \\ 
			& (4.248) \\ 
			& \\ 
			Income x Professional & $-$0.002$^{***}$ \\ 
			& (0.001) \\ 
			& \\ 
			Constant & 21.142$^{***}$ \\ 
			& (2.804) \\ 
			& \\ 
			\hline \\[-1.8ex] 
			Observations & 98 \\ 
			R$^{2}$ & 0.787 \\ 
			Adjusted R$^{2}$ & 0.780 \\ 
			Residual Std. Error & 8.012 (df = 94) \\ 
			F Statistic & 115.878$^{***}$ (df = 3; 94) \\ 
			\hline 
			\hline \\[-1.8ex] 
			\textit{Note:}  & \multicolumn{1}{r}{$^{*}$p$<$0.1; $^{**}$p$<$0.05; $^{***}$p$<$0.01} \\ 
		\end{tabular} 
	\end{table} 
	\item [(c)]
	Write the prediction equation based on the result.
	
		\begin{equation}
		\hat{y} = \hat{\beta}_0 + \hat{\beta}_1 X_i + \hat{\beta}_2 D_i + \hat{\beta}_3(X_i\times D_i)
	\end{equation}
	\begin{equation}
		\widehat{prestige} = 21.142  + 0.003\text{(income)} + 37.781\text{(professional)} - 0.002\text{(income}\times\text{professional)}
	\end{equation}
	
	\item [(d)]
	Interpret the coefficient for \texttt{income}.
	
	\noindent The coefficient for income tells us the conditional effect of income on prestigious jobs for blue collar workers. This means that for every unit increase in income for blue or white collar workers, on average there is a 0.003 increase in job prestige. This shows that there is an association between higher income and more prestigious jobs among non-professionals. This is significant at the 0.001 level. \\
		
	\item [(e)]
	Interpret the coefficient for \texttt{professional}.
	
	\noindent The coefficient for professional tells us the difference in prestige when having a professional job compared to a blue or white collar job when income is at 0. This means that moving from a blue or white colour jobs to a professional job is associated with a 37.781 increase in job prestige when income is 0.    \\
	
	\item [(f)]
	What is the effect of a \$1,000 increase in income on prestige score for professional occupations? In other words, we are interested in the marginal effect of income when the variable \texttt{professional} takes the value of $1$. Calculate the change in $\hat{y}$ associated with a \$1,000 increase in income based on your answer for (c).
	
	\noindent To find the change in $\hat{y}$ when income increases by 1000 I calculated the difference in prestige score between an income of a professional earning 0 and a professional earning 1000.	$\hat{\beta}_0$ and $\hat{\beta}_2 D_i$ cancel out when calculating this difference, leaving the following equation. I also double checked the calculation with the full prediction equations in R, getting the same result. The difference in prestige score is 0.8452. \\
	\begin{equation}
		\widehat{prestige} = 0.003\text{(income)} - 0.002\text{(income}\times\text{professional)}
	\end{equation}
	\begin{equation}
		\widehat{prestige} = 0.003{(1000)} - 0.002{(1000\times1)}
	\end{equation}
	\begin{equation}
		\widehat{prestige} = (0.003 - 0.002)(1000)
	\end{equation}
	\begin{equation}
		\widehat{prestige} = 0.8452
	\end{equation}
	\lstinputlisting[language=R, firstline=61,lastline=62]{PS04.my.answers.el.la.R}
	\begin{lstlisting}[language=R]
		0.8452	\end{lstlisting}
	\lstinputlisting[language=R, firstline=65,lastline=65]{PS04.my.answers.el.la.R}
	\begin{lstlisting}[language=R]
	0.8452	\end{lstlisting}
	
	\item [(g)]
	What is the effect of changing one's occupations from non-professional to professional when her income is \$6,000? We are interested in the marginal effect of professional jobs when the variable \texttt{income} takes the value of $6,000$. Calculate the change in $\hat{y}$ based on your answer for (c).
	
	\noindent To calculate the change in $\hat{y}$ between a professional and non-professional with an income of 6000 I first removed $\hat{\beta}_0$ and $\hat{\beta}_1  X$ as they would cancel out when calculating the change. The value for a non-professional is 0, so to calculate the difference I inputted 1 for professional and 6000 for income, finding a change of 23.827.  \\
	\newpage
	\begin{equation}
		\widehat{prestige} = 31.781\text{(professional)} - 0.002\text{(income}\times\text{professional)}
	\end{equation}
	\begin{equation}
		\widehat{prestige} = 37.781{(1)} - 0.002{(6000\times1)}
	\end{equation}
	\begin{equation}
		\widehat{prestige} = 37.781 - 0.002{(6000)}
	\end{equation}
	\begin{equation}
		\widehat{prestige} = 23.82708
	\end{equation}
	\lstinputlisting[language=R, firstline=69,lastline=70]{PS04.my.answers.el.la.R}
	\begin{lstlisting}[language=R]
	23.82708	\end{lstlisting}
	\lstinputlisting[language=R, firstline=72,lastline=72]{PS04.my.answers.el.la.R}
	\begin{lstlisting}[language=R]
	23.82708	\end{lstlisting}
	
	
\end{enumerate}


\section*{Question 2: Political Science}
\vspace{.25cm}
\noindent 	Researchers are interested in learning the effect of all of those yard signs on voting preferences.\footnote{Donald P. Green, Jonathan	S. Krasno, Alexander Coppock, Benjamin D. Farrer,	Brandon Lenoir, Joshua N. Zingher. 2016. ``The effects of lawn signs on vote outcomes: Results from four randomized field experiments.'' Electoral Studies 41: 143-150. } Working with a campaign in Fairfax County, Virginia, 131 precincts were randomly divided into a treatment and control group. In 30 precincts, signs were posted around the precinct that read, ``For Sale: Terry McAuliffe. Don't Sellout Virgina on November 5.'' \\

Below is the result of a regression with two variables and a constant.  The dependent variable is the proportion of the vote that went to McAuliff's opponent Ken Cuccinelli. The first variable indicates whether a precinct was randomly assigned to have the sign against McAuliffe posted. The second variable indicates
a precinct that was adjacent to a precinct in the treatment group (since people in those precincts might be exposed to the signs).  \\

\vspace{.5cm}
\begin{table}[!htbp]
	\centering 
	\textbf{Impact of lawn signs on vote share}\\
	\begin{tabular}{@{\extracolsep{5pt}}lccc} 
		\\[-1.8ex] 
		\hline \\[-1.8ex]
		Precinct assigned lawn signs  (n=30)  & 0.042\\
		& (0.016) \\
		Precinct adjacent to lawn signs (n=76) & 0.042 \\
		&  (0.013) \\
		Constant  & 0.302\\
		& (0.011)
		\\
		\hline \\
	\end{tabular}\\
	\footnotesize{\textit{Notes:} $R^2$=0.094, N=131}
\end{table}

\begin{enumerate}
	\item [(a)] Use the results from a linear regression to determine whether having these yard signs in a precinct affects vote share (e.g., conduct a hypothesis test with $\alpha = .05$).
	
	\noindent $H_0: \beta_1 = 0 $ \\
	\noindent $H_1: \beta_1 \neq 0 $ \\
	
	\noindent Null hypothesis: There is no linear relationship between having yard signs against opponent in a precinct and vote share. \\
	\noindent Alternative hypothesis: There is a linear relationship between having yard signs against opponent in a precinct and vote share. \\
	
	\noindent I first calculated a test statistic by diving the coefficient for having a yard sign in a precinct by its standard error. I then calculated the p-value with the test statistic. The degrees of freedom were 128 (131-2-1). \\
	
	\lstinputlisting[language=R, firstline=77,lastline=78]{PS04.my.answers.el.la.R}
	
	\begin{lstlisting}[language=R]
		2.625	\end{lstlisting}
	
	\lstinputlisting[language=R, firstline=80,lastline=80]{PS04.my.answers.el.la.R}
	
	\begin{lstlisting}[language=R]
	0.009745963	\end{lstlisting}
	
	\noindent The test statistic was 2.625 and the p-value was 0.0097. As the p-value 0.0097 $<$ 0.05, we reject the null hypothesis that there is no linear relationship between having yard signs against opponent in a precinct and vote share.  \\
		
	\item [(b)]  Use the results to determine whether being
	next to precincts with these yard signs affects vote
	share (e.g., conduct a hypothesis test with $\alpha = .05$).
	
	\noindent $H_0: \beta_1 = 0 $ \\
	\noindent $H_1: \beta_1 \neq 0 $ \\
	
	\noindent Null hypothesis: There is no linear relationship between having yard signs against opponent next to a precinct and vote share. \\
	\noindent Alternative hypothesis: There is a linear relationship between having yard signs against opponent next to a precinct and vote share. \\
	
	\noindent I first calculated a test statistic by diving the coefficient for having a yard sign next to a precinct by its standard error. I then calculated the p-value with the test statistic. The degrees of freedom were 128 (131-2-1). \\
	
	\lstinputlisting[language=R, firstline=85,lastline=86]{PS04.my.answers.el.la.R}
	
	\begin{lstlisting}[language=R]
		3.230769	\end{lstlisting}
	
	\lstinputlisting[language=R, firstline=88,lastline=89]{PS04.my.answers.el.la.R}
	
	\begin{lstlisting}[language=R]
		0.00156946	\end{lstlisting}
	
	\noindent The test statistic was 3.231 and the p-value was 0.0016. As the p-value 0.0016 $<$ 0.05, we reject the null hypothesis that there is no linear relationship between having yard signs against opponent in an adjacent precinct and vote share. \\
	
	\item [(c)] Interpret the coefficient for the constant term substantively.
	
	\noindent The coefficient for the constant term represents the expected vote share for Ken Cuccinelli, 0.302, in precincts with no yard signs against his opponent which are also not adjacent to a precinct with yard signs against his opponent. Precincts with neither the treatment which were not adjacent to it are expected to have 0.042 or 4.2\% less of the vote share. \\
	
	\item [(d)] Evaluate the model fit for this regression.  What does this	tell us about the importance of yard signs versus other factors that are not modeled?
	
	To assess the model fit of this regression we can look at the ${R}^2$ term of 0.094. This means that 9.4\% of the variation in vote share can be explained by the variation in treatment groups in precincts (vote signs inside of, adjacent to or neither). Despite the effect of having or being adjacent to these signs on vote share is significant, a large proportion of variation is still unexplained by the model. This indicates that there are other factors which also contribute to vote share which are not included in the model.
	
\end{enumerate}  


\end{document}
